\documentclass[11pt]{article}
\usepackage[utf8]{inputenc}
\usepackage[margin=1in]{geometry}
\usepackage{amsmath}

\title{Sketched Spectral Embeddings with Gap-Dependent Guarantees}
\author{Razz \and Lena Fischer}
\date{Extended abstract, Workshop on Graph Learning 2025}

\begin{document}
\maketitle

\section*{Summary}

Spectral embeddings of large graphs are usually computed with iterative eigensolvers whose cost grows quickly with the number of eigenvectors requested. We study a randomized alternative: project the graph Laplacian onto the range of a small Gaussian sketch, refined by a few power iterations, and solve the resulting small eigenproblem exactly.

Our main result bounds the distance between the sketched embedding and the exact one by $C \sqrt{n / p} \, \gamma^{q}$. Here $p$ is the oversampling parameter, $q$ the number of power iterations, and $\gamma < 1$ the ratio of consecutive shifted eigenvalues. The bound shows that the number of iterations needed depends on the spectral gap and only logarithmically on the size of the graph.

On three citation networks the sketched embedding reaches the accuracy of the exact one with four power iterations and a tenfold reduction in memory. A full version of this work, with proofs and experiments on larger graphs, is in preparation.

\end{document}
